% HD-S39. GPT-6 Astra Pro, /root/sixfolds; inherited effort not exposed.
% 21 November 2026. Exact reflection-free extension to imaginary quadratic fields.
\subsection{HD-S39: reflection-free sources over imaginary quadratic fields}
\label{sec:imaginary-reflection-free-extension}

Let \(K\) be any imaginary quadratic field, and use an actual
finite absolutely irreducible matrix group \(G\subset GL_m(K)\)
whose character field is exactly \(K\). Use its specified
primitive rational group-algebra corner in the Jacobian of a
connected full-\(G\) sphere cover. Consider a signature-\((3,3)\)
sixfold period family of rank nine. The following extensions are
proved here, with the coefficient and quotient maps retained.

Every reflection-free coefficient representation with
\(9\le m\le14\) is complex imprimitive and has no fixed
nontrivial tensor-product decomposition. In coefficient rank
nine, no such reflection-free source exists when
\(\mu(K)=\{\pm1\}\). In ranks ten through fourteen, a
reflection-free \(K\)-primitive source normalizes no proper
nontrivial unital semisimple \(K\)-subalgebra of
\(\operatorname{End}_K(V)\).
The latter statement in particular excludes every
\(K\)-primitive source with an actual Galois-stable complex
block algebra; it does not assert the existence of such an algebra
for every complex imprimitive representation.

\paragraph{1. The primitive and tensor exclusions are independent of the field.}
The Euler-characteristic and branch-dimension formula (S37.1)
and all codimension inequalities in HD-S37 use only the actual
complex dimensions. The projective tensor-factor homomorphisms,
their cyclic-quotient invariant subspaces, and the actual
scalar/permutation maps (S37.3)--(S37.6) are unchanged.
The contradiction for the deleted permutation representation is
that its full character is real. For every imaginary quadratic
field \(K\), one has \(K\cap\mathbb R=\mathbb Q\), so
this again contradicts the exact character field \(K\).
The exceptional degree-ten normalizer in HD-S37 preserves its
nondegenerate alternating form because the inertia fixed spaces
have dimension at least six; that also makes the full character
real and gives the same contradiction. This argument uses the
complex form of the actual classified module and does not assert
that its displayed Gaussian matrix realization is a model over
the new field. Thus the specified published GS classification
and all its required hypotheses apply exactly as before.

\paragraph{2. Rank nine with only two roots of unity.}
In degree nine, all twelve inertia elements are actual
bireflections. Theorem~8.4 of Guralnick--Saxl gives a line system,
normal diagonal subgroup \(T_0\), and natural permutation
quotient \(A_9\) or \(S_9\), exactly as (S37.8).
The proof that \(T_0\) cannot be scalar remains literal.
The fixed-cycle holonomy calculation and the trace-five
centralizer argument exclude a transposition permutation.
All remaining inertia permutations are 3-cycles or double
transpositions and have determinant one. A scalar in \(T_0\)
now satisfies
\[
                    \zeta\in\{\pm1\},\qquad\zeta^9=1,
                    \qquad\text{hence }\zeta=1.             \tag{S39.1}
\]
The group would consequently be \(A_9\), and its line
stabilizer \(A_8\) has trivial one-dimensional character.
The actual map \(g\otimes1\mapsto\rho(g)v\) in (S37.9)
would identify the representation with the reducible permutation
module. This is the same contradiction, with (S39.1) replacing
the Gaussian fourth-root calculation.

For non-scalar \(T_0\), 2-transitivity makes its nine line
characters distinct. The cyclotomic Galois action fixing \(K\)
commutes with the \(A_9\) or \(S_9\) permutation action,
whose centralizer is trivial. Hence each character is \(K\)-valued.
The literal projectors
\[
       p_j=|T_0|^{-1}\sum_{t\in T_0}\chi_j(t)^{-1}\rho(t)
             \in M_9(K)                                    \tag{S39.2}
\]
give nine actual rank-one \(K\)-subspaces. Choose one, with
stabilizer \(H\), nonzero vector \(v\), and line character
\(\chi:H\to K^*\). The coset map
\[
       e_j\longmapsto\rho(t_j)v
       \quad\text{identifies }V\text{ with }
       \operatorname{Ind}_H^G\chi                            \tag{S39.3}
\]
over \(K\). Since \(H\) is finite,
\(\chi(H)\subset\mu(K)=\{\pm1\}\).
Every induced character value is therefore an integer sum of
\(\pm1\). Its character field is rational, contradicting
the stipulated field \(K\). This proves nonexistence of the
specified rank-nine reflection-free source.

In particular this applies to
\[
 \mathbb Q(\sqrt{-2}),\ \mathbb Q(\sqrt{-7}),\
 \mathbb Q(\sqrt{-11}),\ \mathbb Q(\sqrt{-19}),\
 \mathbb Q(\sqrt{-43}),\ \mathbb Q(\sqrt{-67}),\
 \mathbb Q(\sqrt{-163}).                                   \tag{S39.4}
\]
The roots-of-unity assertion used here follows directly from
\([\mathbb Q(\zeta_n):\mathbb Q]\le2\): the possible
orders are \(1,2,3,4,6\), and the latter three nontrivial
possibilities give \(\mathbb Q(\sqrt{-3})\) or
\(\mathbb Q(i)\). No class-number calculation is needed
for this source restriction. The scalar-diagonal argument for
the Eisenstein field has not been included: its third roots
do satisfy \(\zeta^9=1\), so (S39.1) does not apply there.

\paragraph{3. Normalized algebras in ranks ten through fourteen.}
For a \(K\)-primitive source, the algebraic reduction in
paragraphs 2--3 of HD-S38 uses only characteristic zero and the
actual normalization map. Product centres would give proper
\(K\)-blocks; centre \(K\) gives the explicit fixed tensor
map via matrix units; a proper centre field \(E/K\) gives
the actual homomorphism \(G\to\operatorname{Aut}_K(E)\).
Transitivity on embeddings makes the latter a regular Galois
quotient. The codimension and product-one inequalities are
identical, so the sole numerical endpoint is again
\[
       m=14,\quad[E:K]=2,\quad \dim_E V=7,\quad b=12,
                                                               \tag{S39.5}
\]
with two outside involutions and ten \(E\)-reflections.
Their normal closure is exactly the field-action kernel \(H\).
Its seven-dimensional complex representation is irreducible and
its character is not contained in \(K\), by the explicit
\(W\oplus tW\) and \(U\otimes\mathbb C^2\) maps of HD-S38.

The same rank-seven reflection classification therefore gives
\(H=G(d,e,7)\) in its natural representation; all real
rank-seven possibilities have rational character. The two
diagonal trace elements give the actual identity (S38.10),
so \(\zeta_d\in E\). Since \([E:\mathbb Q]=4\),
the elementary cyclotomic degree bound first gives
\[
                    d\in\{1,2,3,4,5,6,8,10,12\}.           \tag{S39.6}
\]
The field \(\mathbb Q(\zeta_5)=\mathbb Q(\zeta_{10})\)
has cyclic Galois group of order four; its unique quadratic
subfield is fixed by complex conjugation and is real. It
cannot equal a field containing the imaginary quadratic \(K\).
This excludes \(d=5,10\). The cases \(d=1,2\), and any
case with \(\zeta_d\in K\), have character in \(K\)
and were already excluded. The remaining list is contained in
\[
                             \{3,4,6,8,12\}.                \tag{S39.7}
\]
For every retained case \(E=K(\zeta_d)\), so its
nontrivial automorphism fixes the coordinate axes and maps
\(G(d,e,7)\) onto itself by a coprime phase power.
The finite intertwiner system is solved over \(E\), giving
the actual restricted-scalar \(K\)-map. The normalizer
argument of HD-S38 works for every \(d\) in (S39.7),
including \(d=4\): consecutive phases have normalized
inner product \(\cos(\pi/4)>1/2\), so the same root-line
graph recovers the coordinate pair-planes. The \(d=3\)
orthogonality-neighborhood argument is unchanged.
It follows that the full semilinear group preserves seven
\(E\)-coordinate lines, which are actual two-dimensional
\(K\)-blocks, contradicting \(K\)-primitivity.

For complete coordinate precision, choose any
\(\beta\in E\setminus K\) with \(\operatorname{Tr}_{E/K}\beta=0\)
and put \(q=\beta^2\in K\). The restriction-of-scalars and
field-action matrices are literally
\[
 u+v\beta\longmapsto\begin{pmatrix}u&qv\\v&u\end{pmatrix},
 \qquad\tau\longmapsto\begin{pmatrix}1&0\\0&-1\end{pmatrix}.
                                                               \tag{S39.8}
\]
Thus the new-field argument retains the coefficient maps rather
than identifying different field actions abstractly.

The actual \(K\)-block reductions, when they occur, continue
through the same quotient curves and literal primitive idempotent
equalities (S38.2)--(S38.3). Their pullback polarization multiplier
is \(|\ker\sigma|\), with determinant norm
\(|\ker\sigma|^6\). Any further determinant conclusion for
those smaller coefficient groups uses their separately established
source analysis; none is assumed merely from the dimension of a
parameter space.
